\documentclass[a4paper,12pt]{article}

\usepackage[T2A]{fontenc}
\usepackage[utf8]{inputenc}
\usepackage[russian]{babel}

\usepackage{amsmath, amssymb, amsthm, mathtools}
\usepackage{helvet}
\renewcommand{\familydefault}{\sfdefault}
\usepackage{icomma}
\usepackage[left=18mm,right=18mm,top=18mm,bottom=20mm]{geometry}
\usepackage{microtype}
\usepackage{tikz}
\usepackage{tkz-euclide}
\usepackage{pgfplots}
\pgfplotsset{compat=1.18}
\usepackage{tabularx, booktabs, longtable}
\usepackage{enumitem}
\usepackage{hyperref}
\usepackage{parskip}
\usepackage{fancyhdr}
\usepackage{setspace}
\usepackage{xcolor}

\definecolor{accent}{HTML}{008080}
\definecolor{darkaccent}{HTML}{004D4D}
\definecolor{textgray}{HTML}{333333}
\definecolor{softgray}{HTML}{666666}
\definecolor{lightaccent}{HTML}{E8F4F4}

\geometry{a4paper}
\onehalfspacing
\hypersetup{
    colorlinks=true,
    linkcolor=darkaccent,
    urlcolor=darkaccent
}

\pagestyle{fancy}
\fancyhf{}
\fancyhead[L]{\small Чек-лист: прикладные задачи ЕГЭ}
\fancyhead[R]{\small Ксения Клыкова}
\fancyfoot[C]{\small \thepage}
\renewcommand{\headrulewidth}{0.4pt}

\setlength{\parindent}{0pt}
\setlength{\parskip}{6pt}

\setlist[itemize]{leftmargin=1.2cm, itemsep=2pt, topsep=3pt}
\setlist[enumerate]{leftmargin=1.2cm, itemsep=2pt, topsep=3pt}

\newcommand{\sectiontitle}[1]{\section*{\textcolor{darkaccent}{#1}}\addcontentsline{toc}{section}{#1}}
\newcommand{\subsectiontitle}[1]{\subsection*{\textcolor{accent}{#1}}}
\newcommand{\important}[1]{\textbf{\textcolor{darkaccent}{#1}}}

\begin{document}

\begin{titlepage}
\thispagestyle{empty}
\begin{center}
\vspace*{28mm}

{\Large \textcolor{softgray}{Чек-лист}}

\vspace{9mm}

{\Huge \textbf{Прикладные задачи ЕГЭ}}

\vspace{6mm}

{\Large \textbf{Формулы, степени, дроби и выбор физически верного ответа}}

\vspace{22mm}

{\large Лёвин Артём Александрович эксклюзивно для Ксении Клыковой}

\vspace{12mm}

{\large \today}

\vfill

\begin{tikzpicture}
\draw[accent, line width=1.1pt]
  (0,0) .. controls (2,0.85) and (4,-0.85) .. (6,0)
        .. controls (8,0.85) and (10,-0.85) .. (12,0);
\end{tikzpicture}

\vspace{14mm}
\end{center}
\end{titlepage}

\tableofcontents
\newpage

\sectiontitle{1. Главная идея прикладных задач}

Прикладная задача в ЕГЭ часто выглядит как задача по физике, технике или экономике, но математически обычно сводится к аккуратной работе с формулой.

\important{Основной сценарий решения:}

\begin{enumerate}
\item Выписать рабочую формулу.
\item Понять, какую величину нужно найти.
\item Перевести текстовое условие в математическое равенство.
\item Подставить данные.
\item Аккуратно упростить выражение.
\item Решить полученное уравнение.
\item Выбрать ответ, подходящий по смыслу задачи.
\end{enumerate}

\important{Ключевой принцип:} в таких задачах чаще всего ошибаются не в сложной математике, а в арифметике, степенях, дробях и переводе текста условия в уравнение.

\sectiontitle{2. Перевод текста задачи в математику}

\subsectiontitle{Четверть первоначального объёма}

Если в задаче бак имеет постоянную форму, то объём воды пропорционален высоте уровня воды. Поэтому если осталась четверть первоначального объёма, то осталась четверть первоначальной высоты:
\[
h=\frac{h_0}{4}.
\]

Например, если
\[
h_0=20,
\]
то
\[
h=\frac{20}{4}=5.
\]

\important{Важно:} такой переход допустим, когда площадь горизонтального сечения бака постоянна или в условии задачи подразумевается именно такая модель.

\subsectiontitle{Слова, на которые нужно реагировать}

\begin{center}
\renewcommand{\arraystretch}{1.35}
\begin{tabularx}{\linewidth}{>{\bfseries}p{0.32\linewidth} X}
\toprule
Фраза в условии & Математический смысл \\
\midrule
Осталась четверть первоначального объёма & \(h=\dfrac{h_0}{4}\), если объём пропорционален высоте \\
Не менее \(A\) & Обычно означает \(\ge A\), но в задачах на точное значение часто берут граничный случай \(=A\) \\
Найти через сколько секунд & Неизвестное чаще всего \(t\) \\
Найти линейный размер & Неизвестное чаще всего \(L\) \\
Найти температуру & Неизвестное чаще всего \(T\) \\
\bottomrule
\end{tabularx}
\end{center}

\sectiontitle{3. Задача на вытекание воды}

\subsectiontitle{Рабочая формула}

В задачах такого типа высота уровня воды может задаваться формулой:
\[
h=h_0-\frac{\sqrt{2gh_0}}{k}\,t+\frac{g}{2k^2}\,t^2.
\]

Здесь:
\begin{itemize}
\item \(h_0\) — начальная высота уровня воды;
\item \(h\) — высота уровня воды через \(t\) секунд;
\item \(g\) — ускорение свободного падения;
\item \(k\) — постоянная из условия задачи;
\item \(t\) — время.
\end{itemize}

\subsectiontitle{Типовой пример}

Пусть
\[
h_0=20,\qquad g=10,\qquad k=500.
\]
Нужно найти время, когда в баке останется четверть первоначального объёма.

Так как осталась четверть первоначального объёма, берём
\[
h=\frac{20}{4}=5.
\]

Подставляем:
\[
5=20-\frac{\sqrt{2\cdot 10\cdot 20}}{500}\,t+\frac{10}{2\cdot 500^2}\,t^2.
\]

Упрощаем:
\[
\sqrt{2\cdot 10\cdot 20}=\sqrt{400}=20,
\]
поэтому
\[
5=20-\frac{20}{500}t+\frac{10}{2\cdot 500^2}t^2.
\]

Так как
\[
\frac{20}{500}=\frac{1}{25},\qquad
\frac{10}{2\cdot 500^2}=\frac{5}{250000}=\frac{1}{50000},
\]
получаем
\[
5=20-\frac{t}{25}+\frac{t^2}{50000}.
\]

Переносим всё в одну сторону:
\[
0=15-\frac{t}{25}+\frac{t^2}{50000}.
\]

Домножаем на \(50000\):
\[
0=750000-2000t+t^2.
\]

То есть
\[
t^2-2000t+750000=0.
\]

Дискриминант:
\[
D=2000^2-4\cdot 750000=4000000-3000000=1000000.
\]
\[
\sqrt{D}=1000.
\]

Корни:
\[
t_1=\frac{2000-1000}{2}=500,
\qquad
t_2=\frac{2000+1000}{2}=1500.
\]

\important{По смыслу задачи подходит \(t=500\).} Если через \(500\) секунд осталась четверть воды, то второй корень \(1500\) секунд относится уже к продолжению математической модели после физически важного момента и в ответ не берётся.

\subsectiontitle{Что обязательно проверить}

\begin{itemize}
\item Верно ли Вы перевели «четверть объёма» в \(h=\dfrac{h_0}{4}\).
\item Не забыли ли скобки в выражении \(2k^2\).
\item Верно ли сократили дроби.
\item Не потеряли ли множитель \(t\) или \(t^2\).
\item После решения квадратного уравнения выбрали ли ответ по смыслу задачи.
\end{itemize}

\sectiontitle{4. Задачи с кубом величины}

\subsectiontitle{Рабочая формула}

В прикладных задачах может встретиться формула вида:
\[
T=\rho gL^3-mg.
\]

Здесь:
\begin{itemize}
\item \(T\) — сила натяжения троса;
\item \(\rho\) — плотность;
\item \(g\) — ускорение свободного падения;
\item \(L\) — линейный размер аппарата;
\item \(m\) — масса.
\end{itemize}

Если нужно найти \(L\), сначала выражают \(L^3\), а затем извлекают кубический корень.

\subsectiontitle{Типовой пример}

Пусть
\[
T=1000,\qquad \rho=1000,\qquad g=10,\qquad m=25.
\]

Подставляем:
\[
1000=1000\cdot 10\cdot L^3-25\cdot 10.
\]

Упрощаем:
\[
1000=10000L^3-250.
\]

Переносим:
\[
1250=10000L^3.
\]

Тогда
\[
L^3=\frac{1250}{10000}=\frac{125}{1000}=\frac{1}{8}.
\]

Значит,
\[
L=\sqrt[3]{\frac{1}{8}}=\frac{1}{2}=0{,}5.
\]

\subsectiontitle{Как извлекать корень через степени}

Если
\[
L^3=\frac{125}{1000},
\]
то можно записать:
\[
L=\left(\frac{125}{1000}\right)^{\frac13}.
\]

Представляем числитель и знаменатель как кубы:
\[
125=5^3,\qquad 1000=10^3.
\]

Тогда
\[
L=\left(\frac{5^3}{10^3}\right)^{\frac13}
=\frac{5^{3\cdot \frac13}}{10^{3\cdot \frac13}}
=\frac{5}{10}=0{,}5.
\]

\important{Главная идея:} чтобы из \(L^3\) получить \(L\), нужно возвести обе части равенства в степень \(\dfrac13\).

\sectiontitle{5. Степени: как убирать показатель}

\subsectiontitle{Общее правило}

Если неизвестная стоит в степени, то нужно применить обратную степень.

\[
x^3=a \quad \Longrightarrow \quad x=a^{\frac13}.
\]

\[
x^4=a \quad \Longrightarrow \quad x=a^{\frac14}.
\]

\[
x^{\frac43}=a \quad \Longrightarrow \quad x=a^{\frac34}.
\]

\subsectiontitle{Почему так работает}

Если степень возводится в степень, показатели перемножаются:
\[
\left(x^a\right)^b=x^{ab}.
\]

Поэтому:
\[
\left(x^3\right)^{\frac13}=x^{3\cdot \frac13}=x.
\]

И так же:
\[
\left(x^{\frac43}\right)^{\frac34}=x^{\frac43\cdot\frac34}=x.
\]

\subsectiontitle{Пример с дробной степенью}

Пусть
\[
t^{\frac43}=\frac{16}{25}.
\]

Чтобы получить \(t\), возводим обе части в степень \(\dfrac34\):
\[
t=\left(\frac{16}{25}\right)^{\frac34}.
\]

Представляем:
\[
16=2^4,\qquad 25=5^2.
\]

Тогда
\[
t=\left(\frac{2^4}{5^2}\right)^{\frac34}
=\frac{2^3}{5^{\frac32}}.
\]

Это выражение не является красивым целым или десятичным числом. В экзаменационной задаче такой результат может быть допустим, если ответ требуется в таком виде, но в прикладных задачах ЕГЭ чаще данные подобраны так, чтобы после преобразований получался простой ответ.

\sectiontitle{6. Задачи со степенями десяти}

\subsectiontitle{Как работать с научной записью}

Научная запись имеет вид:
\[
a\cdot 10^n.
\]

Например:
\[
5{,}7\cdot 10^{-8}=57\cdot 10^{-9}.
\]

Почему так:
\[
5{,}7=57\cdot 10^{-1},
\]
значит
\[
5{,}7\cdot 10^{-8}=57\cdot 10^{-1}\cdot 10^{-8}=57\cdot 10^{-9}.
\]

\important{Правило:} если запятую перенесли вправо на один знак, степень десяти уменьшается на \(1\).

Если запятую перенесли вправо на четыре знака, степень десяти уменьшается на \(4\).

\subsectiontitle{Типичная ловушка}

Например:
\[
1{,}5625\cdot 10^{25}.
\]

Если перенести запятую вправо на четыре знака:
\[
1{,}5625=15625\cdot 10^{-4}.
\]

Тогда
\[
1{,}5625\cdot 10^{25}=15625\cdot 10^{21}.
\]

\important{Ошибка:} записать \(15625\cdot 10^{20}\). Это означает, что степень уменьшили на \(5\), а не на \(4\).

\subsectiontitle{Сигнал ошибки}

Если в задаче данные явно подобраны под красивый ответ, но при извлечении корня появляются неудобные дробные показатели, нужно проверить:

\begin{itemize}
\item перенос запятой;
\item степень числа \(10\);
\item сокращение одинаковых множителей;
\item разложение составных чисел на простые множители.
\end{itemize}

\sectiontitle{7. Формула с четвёртой степенью}

\subsectiontitle{Рабочая формула}

В задачах на излучение может встретиться формула:
\[
P=\sigma ST^4.
\]

Здесь:
\begin{itemize}
\item \(P\) — мощность излучения;
\item \(\sigma\) — постоянная из условия;
\item \(S\) — площадь;
\item \(T\) — температура.
\end{itemize}

Если нужно найти \(T\), сначала выражают \(T^4\), затем берут четвёртый корень:
\[
T=\sqrt[4]{\frac{P}{\sigma S}}.
\]

\subsectiontitle{Типовой порядок действий}

\begin{enumerate}
\item Аккуратно переписать все числа в степенной форме.
\item Подставить в формулу \(P=\sigma ST^4\).
\item Выразить:
\[
T^4=\frac{P}{\sigma S}.
\]
\item Сократить одинаковые степени десяти.
\item Разложить большие числа на простые множители.
\item Представить правую часть как четвёртую степень.
\item Извлечь четвёртый корень.
\end{enumerate}

\subsectiontitle{Пример разложения}

Если после сокращений получилось:
\[
T^4=5^{16}\cdot 2^{12},
\]
то
\[
T=\left(5^{16}\cdot 2^{12}\right)^{\frac14}
=5^4\cdot 2^3.
\]

Считаем:
\[
5^4=625,\qquad 2^3=8.
\]

Тогда
\[
T=625\cdot 8=5000.
\]

\sectiontitle{8. Как быстро раскладывать большие числа}

Большие числа в прикладных задачах часто специально подобраны так, чтобы хорошо сокращаться.

\begin{center}
\renewcommand{\arraystretch}{1.35}
\begin{tabularx}{\linewidth}{>{\bfseries}c X}
\toprule
Число & Разложение \\
\midrule
\(125\) & \(5^3\) \\
\(343\) & \(7^3\) \\
\(625\) & \(5^4\) \\
\(1000\) & \(10^3=2^3\cdot 5^3\) \\
\(15625\) & \(5^6\) \\
\(228\) & \(2^2\cdot 3\cdot 19\) \\
\(57\) & \(3\cdot 19\) \\
\(10000\) & \(10^4=2^4\cdot 5^4\) \\
\bottomrule
\end{tabularx}
\end{center}

\important{Приём:} если числа большие, не перемножайте их сразу. Сначала разложите на множители и сократите.

\sectiontitle{9. Типичные ошибки}

\begin{center}
\renewcommand{\arraystretch}{1.35}
\begin{tabularx}{\linewidth}{>{\bfseries}p{0.34\linewidth} X}
\toprule
Ошибка & Как правильно \\
\midrule
Подставить в формулу все данные, но не перевести текст условия & Сначала понять, чему равна искомая величина: например, \(h=\dfrac{h_0}{4}\). \\
Потерять скобки в \(2k^2\) & Запись \(\dfrac{g}{2k^2}\) означает деление на всё выражение \(2k^2\). \\
Сократить дроби слишком быстро и ошибиться & Лучше сокращать по шагам и переписывать промежуточные выражения. \\
Забыть физический смысл корней & Если квадратное уравнение даёт два корня, нужно выбрать тот, который подходит по условию. \\
Неверно перенести запятую в научной записи & При переносе запятой вправо степень десяти уменьшается. \\
Не разложить большие числа на множители & В задачах со степенями это почти всегда замедляет решение и повышает риск ошибки. \\
Неверно убрать степень у неизвестной & Для \(x^3\) нужна степень \(\dfrac13\), для \(x^4\) — степень \(\dfrac14\), для \(x^{4/3}\) — степень \(\dfrac34\). \\
\bottomrule
\end{tabularx}
\end{center}

\sectiontitle{10. Мини-чек-лист перед ответом}

Перед записью ответа проверьте:

\begin{itemize}
\item Все ли данные подставлены в нужные места формулы?
\item Не забыта ли величина, скрытая в тексте условия?
\item Нет ли ошибки в степени числа \(10\)?
\item Не потерян ли квадрат, куб или четвёртая степень?
\item Правильно ли выбран корень по смыслу задачи?
\item Ответ записан в той величине, которую спрашивали?
\end{itemize}

\newpage

\sectiontitle{Тренировочный блок}

\textbf{Средний уровень}

\textbf{1.} Высота уровня воды в баке задаётся формулой
\[
h=h_0-\frac{\sqrt{2gh_0}}{k}\,t+\frac{g}{2k^2}\,t^2.
\]
Дано:
\[
h_0=20,\qquad g=10,\qquad k=400.
\]
Через сколько секунд в баке останется четверть первоначального объёма? В ответ запишите меньший физически подходящий корень.

\textbf{2.} Сила натяжения троса вычисляется по формуле
\[
T=\rho gL^3-mg.
\]
Дано:
\[
T=2600,\qquad \rho=1000,\qquad g=10,\qquad m=83.
\]
Найдите \(L\).

\textbf{3.} Представьте число в виде \(a\cdot 10^n\) без десятичной запятой в первом множителе:
\[
1{,}5625\cdot 10^{25}.
\]

\textbf{Выше среднего}

\textbf{4.} Высота уровня воды в баке задаётся формулой
\[
h=h_0-\frac{\sqrt{2gh_0}}{k}\,t+\frac{g}{2k^2}\,t^2.
\]
Дано:
\[
h_0=45,\qquad g=10,\qquad k=300.
\]
Через сколько секунд в баке останется четверть первоначального объёма? В ответ запишите меньший физически подходящий корень.

\textbf{5.} Сила натяжения троса вычисляется по формуле
\[
T=\rho gL^3-mg.
\]
Дано:
\[
T=3290,\qquad \rho=1000,\qquad g=10,\qquad m=140.
\]
Найдите \(L\).

\textbf{6.} Решите уравнение:
\[
x^3=\frac{343}{1000}.
\]

\textbf{7.} Решите уравнение:
\[
y^4=625\cdot 16.
\]
В ответ запишите положительное значение \(y\).

\textbf{8.} В формуле
\[
P=\sigma ST^4
\]
даны:
\[
P=15625\cdot 10^{21},\qquad
\sigma=57\cdot 10^{-9},\qquad
S=\frac{10^{-20}}{228}.
\]
Найдите \(T\).

\textbf{9.} Найдите ошибку в преобразовании:
\[
1{,}5625\cdot 10^{25}=15625\cdot 10^{20}.
\]
Запишите верную форму.

\textbf{Сложный уровень}

\textbf{10.} В формуле
\[
P=\sigma ST^4
\]
даны:
\[
P=78125\cdot 10^{20},\qquad
\sigma=25\cdot 10^{-8},\qquad
S=\frac{10^{-18}}{128}.
\]
Найдите \(T\).

\newpage

\sectiontitle{Ответы}

\begin{center}
\renewcommand{\arraystretch}{1.45}
\begin{tabularx}{\linewidth}{>{\bfseries}c X}
\toprule
№ & Ответ \\
\midrule
1 & \(400\) \\
2 & \(0{,}7\) \\
3 & \(15625\cdot 10^{21}\) \\
4 & \(450\) \\
5 & \(0{,}9\) \\
6 & \(0{,}7\) \\
7 & \(10\) \\
8 & \(5000\) \\
9 & Ошибка: степень десяти уменьшили на \(5\), а нужно на \(4\). Верно: \(1{,}5625\cdot 10^{25}=15625\cdot 10^{21}\). \\
10 & \(10000\) \\
\bottomrule
\end{tabularx}
\end{center}

\end{document}